\begin{solapartat}
\punts{0,25} La força mesurada per la balança és igual a la seva lectura en kg per $g$. Així, la força magnètica és igual a la diferència en la lectura de la balança per $g$:
\[ F = (0,243 - 0,239)\times 9,81 = 0,0392\ \mathrm{N} \]

\punts{0,25} Per altre banda, la força aplicada pel camp magnètic sobre un fil de corrent és:
\[ \vec{F}_m = I\vec{l}\times\vec{B} \]

\punts{0,25} Com $\vec{l}$ i $\vec{B}$ són perpendiculars, el mòdul de la força magnètica és:
\[ F_m = IlB \]

\punts{0,25} I finalment, com de la lectura de la balança sabem que $F_m = 0,0392\ \mathrm{N}$:
\[ B = \frac{F_m}{Il} = \frac{0,0392}{8\ 0,15} = 0,0327\ \mathrm{T} \]
% DUBTE: al denominador la pauta escriu "8 0,15" sense signe de multiplicar (8 × 0,15).

\punts{0,25} Observeu que $\vec{F}_m$ serà perpendicular a $\vec{l}$ i $\vec{B}$ per tant serà vertical. El camp magnètic va dirigit del pol nord cap el pol sud, i tenint en compte el sentit del corrent, si apliquem la regla de la mà dreta tenim que la força magnètica aplicada sobre el fil va dirigida cap amunt. Segons la tercera llei de Newton, la força aplicada per l'imant sobre el fil de corrent tindrà la mateixa magnitud i direcció però sentit oposat al de la força aplicada pel fil de corrent sobre l'imant, per tant, sobre l'imant actua una força dirigida segons l'eix $z$ i sentit cap avall (negatiu) de manera que la balança detectarà un augment de la força aplicada sobre ella.

\figura[0.25]{sol_fig1.png}
\end{solapartat}

\begin{solapartat}
\punts{0,75} Per obtenir la mateixa força hem de crear el mateix camp magnètic:
\[ B = \frac{\mu_0 I}{2\pi r} = 0,0327\ \mathrm{T} \]

\punts{0,5} Per tant,
\[ r = \frac{\mu_0 I}{2\pi B} = \frac{4\pi\times 10^{-7}\times 8}{2\pi\times 0,0327} = 4,90\times 10^{-5}\ \mathrm{m} = 49,0\ \mu\mathrm{m} \]
\end{solapartat}
